A diferencia del paradigma clásico de Django\footnote{\url{https://www.djangoproject.com/}}, en el que el servidor genera y sirve directamente las páginas HTML, SocratiCode adopta una arquitectura API REST desacoplada mediante Django REST Framework (DRF)\footnote{\url{https://www.django-rest-framework.org/}}. En este modelo el \textit{backend} actúa exclusivamente como un servicio de datos, recibe las peticiones HTTP del \textit{frontend}, interactúa con la base de datos PostgreSQL\footnote{\url{https://www.postgresql.org/}} a través del ORM de Django, consultando y persistiendo las entidades, y devuelve las respuestas formateadas en JSON. DRF simplifica este proceso gracias a su sistema de serializadores, que permite convertir automáticamente objetos complejos de Python en estructuras JSON listas para ser consumidas por la interfaz. Para la comunicación asíncrona se utiliza adicionalmente ADRF\footnote{\url{https://github.com/em1208/adrf}}, una extensión que habilita vistas asíncronas nativas sobre DRF, imprescindible para soportar el \textit{streaming} en tiempo real del tutor socrático.